\documentclass[aspectratio=169,10pt]{beamer}

% Shared Beamer setup for Fall 2026 course slides.
\mode<presentation>{
  \usetheme{Madrid}
  \usecolortheme{default}
}

\definecolor{CalPolyGreen}{HTML}{154734}
\definecolor{CalPolyGold}{HTML}{BD8B13}
\definecolor{DeckInk}{HTML}{1F2933}
\definecolor{DeckMuted}{HTML}{5F6B7A}

\setbeamercolor{structure}{fg=CalPolyGreen}
\setbeamercolor{palette primary}{bg=CalPolyGreen,fg=white}
\setbeamercolor{palette secondary}{bg=DeckInk,fg=white}
\setbeamercolor{palette tertiary}{bg=CalPolyGold,fg=DeckInk}
\setbeamercolor{frametitle}{bg=CalPolyGreen,fg=white}
\setbeamercolor{title}{fg=CalPolyGreen}
\setbeamercolor{normal text}{fg=DeckInk,bg=white}
\setbeamercolor{block title}{bg=CalPolyGreen,fg=white}
\setbeamercolor{block body}{bg=black!3,fg=DeckInk}

\setbeamertemplate{navigation symbols}{}
\setbeamertemplate{footline}[frame number]
\setbeamertemplate{itemize item}{\color{CalPolyGold}$\blacktriangleright$}
\setbeamertemplate{itemize subitem}{\color{CalPolyGreen}$\bullet$}

\newcommand{\CourseNumber}{Course}
\newcommand{\CourseTitle}{Course Title}
\newcommand{\LectureNumber}{0}
\newcommand{\LectureTitle}{Lecture Title}
\newcommand{\LectureDate}{Date}
\newcommand{\CourseUnit}{Unit}

\newcommand{\coursenumber}[1]{\renewcommand{\CourseNumber}{#1}}
\newcommand{\coursetitle}[1]{\renewcommand{\CourseTitle}{#1}}
\newcommand{\lecturenumber}[1]{\renewcommand{\LectureNumber}{#1}}
\newcommand{\lecturetitle}[1]{\renewcommand{\LectureTitle}{#1}}
\newcommand{\lecturedate}[1]{\renewcommand{\LectureDate}{#1}}
\newcommand{\courseunit}[1]{\renewcommand{\CourseUnit}{#1}}

\author{Kirk Duran}
\institute{California Polytechnic State University, San Luis Obispo}
\date{\LectureDate}

\newcommand{\makedecktitle}{
  \begin{frame}[plain]
    \vspace{0.25in}
    {\usebeamerfont{title}\color{CalPolyGreen}\LectureTitle\par}
    \vspace{0.18in}
    {\large \CourseNumber{}: \CourseTitle\par}
    \vspace{0.08in}
    {\large Lecture \LectureNumber{} -- \LectureDate\par}
    \vfill
    {\small Kirk Duran \textbar{} Cal Poly, San Luis Obispo\par}
  \end{frame}
}

\newcommand{\placeholdernote}{
  \begin{block}{Placeholder}
    Replace these bullets with the final lecture material, examples, and in-class activities.
  \end{block}
}

\AtBeginSection[]{
  \begin{frame}{Roadmap}
    \tableofcontents[currentsection]
  \end{frame}
}


\graphicspath{{../assets/lecture-04/}}

% If an image file is not yet available, skip the visual slot.
\newcommand{\lectureimage}[2][1.75in]{%
  \IfFileExists{../assets/lecture-04/#2}{%
    \includegraphics[width=\linewidth,height=#1,keepaspectratio]{#2}%
  }{%
  }%
}

\setbeamersize{text margin left=0.42in,text margin right=0.42in}

\coursenumber{CS 1001}
\coursetitle{Introduction to Computing}
\lecturenumber{4}
\lecturetitle{Data Types, Operators, and Conversion}
\lecturedate{August 31, 2026}
\courseunit{1. Computing and Program Execution}

\begin{document}

\makedecktitle

\begin{frame}{Today}

  \begin{itemize}

    \item Why computer systems need data types.

    \item How Python represents integers, floating-point numbers, strings, and Boolean values.

    \item How operand types determine an operator's meaning and result type.

    \item How Python handles compatible and incompatible type combinations.

    \item How explicit conversion produces a value of another type.

  \end{itemize}

  \begin{keyidea}

    Every example remains a self-contained expression; naming and retaining values comes later.

  \end{keyidea}

\end{frame}

\sectionframe{Part 1: Why Data Types Exist}

\begin{frame}{A Group of Bits Does Not Explain Its Own Meaning}

  \begin{columns}[T,onlytextwidth]

    \begin{column}{0.54\textwidth}

      \begin{block}{One possible representation of ``Kirk''}

        \centering
        \texttt{01001011 01101001}\\
        \texttt{01110010 01101011}

      \end{block}

      \begin{itemize}

        \item At the hardware level, text and numbers are represented using bits.

        \item The bit pattern alone does not say whether it represents text, a number, an image, or instructions.

        \item Adding the number \texttt{5} to encoded text has no useful textual meaning.

      \end{itemize}

    \end{column}

    \begin{column}{0.40\textwidth}

      % Search direction: binary bits interpreted as text versus a numerical value.
      \lectureimage[2.22in]{slide4.png}

    \end{column}

  \end{columns}

  \begin{keyidea}

    A representation becomes a meaningful value only when a system knows how to interpret it.

  \end{keyidea}

\end{frame}

\begin{frame}{Different Kinds of Values Need Different Interpretations}

  \begin{columns}[T,onlytextwidth]

    \begin{column}{0.50\textwidth}

      \begin{itemize}

        \item \textbf{Whole numbers:} \texttt{42}, \texttt{-8}

        \item \textbf{Fractional values:} \texttt{3.14}, \texttt{-0.25}

        \item \textbf{Text:} \texttt{"Cal Poly"}

        \item \textbf{Two logical states:} \texttt{True}, \texttt{False}

      \end{itemize}

    \end{column}

    \begin{column}{0.44\textwidth}

      % Search direction: four distinct representations for integer, decimal, text, and two-state data.
      \lectureimage[2.18in]{slide5.png}

    \end{column}

  \end{columns}

  \begin{exampleblockcp}

    Python names these types \texttt{int}, \texttt{float}, \texttt{str}, and \texttt{bool}.

  \end{exampleblockcp}

\end{frame}

\begin{frame}{Every Python Value Has a Type}

  \begin{cpdefinition}

    A \textbf{data type} is a classification that determines how a value is represented, interpreted, and used.

  \end{cpdefinition}

  \begin{columns}[T,onlytextwidth]

    \begin{column}{0.30\textwidth}

      \begin{block}{Representation}

        How Python organizes the information that constitutes the value.

      \end{block}

    \end{column}

    \begin{column}{0.30\textwidth}

      \begin{block}{Operations}

        Which operators and functions are defined for the value.

      \end{block}

    \end{column}

    \begin{column}{0.30\textwidth}

      \begin{block}{Results}

        What value and type an operation produces.

      \end{block}

    \end{column}

  \end{columns}

\end{frame}

\begin{frame}{Representations Introduce Costs and Limits}

  \begin{columns}[T,onlytextwidth]

    \begin{column}{0.53\textwidth}

      \begin{itemize}

        \item A \textbf{fixed-size representation} uses a predetermined amount of storage.

        \item A \textbf{variable-size representation} can grow with the represented value.

        \item Numerical representations may trade exactness for range and efficiency.

        \item A value that does not fit must be shortened, rejected, approximated, or represented differently.

      \end{itemize}

    \end{column}

    \begin{column}{0.41\textwidth}

      % Search direction: fixed-size container contrasted with a growing data container.
      \lectureimage[2.18in]{slide7.png}

    \end{column}

  \end{columns}

\end{frame}

\begin{frame}{A Ten-Character Limit Shaped Feraligatr's Name}

  \begin{columns}[T,onlytextwidth]

    \begin{column}{0.49\textwidth}

      \begin{block}{Ten characters fit}

        \centering
        \texttt{F E R A L I G A T R}\\[0.06in]
        \scriptsize\texttt{1\quad2\quad3\quad4\quad5\quad6\quad7\quad8\quad9\quad10}

      \end{block}

      \begin{block}{Eleven characters do not}

        \centering
        \texttt{F E R A L I G A T O R}

      \end{block}

      \begin{itemize}

        \item Generation II limited English Pok\'emon names to ten characters.

        \item The shortened \emph{Feraligatr} spelling fits that representation.

      \end{itemize}

    \end{column}

    \begin{column}{0.45\textwidth}

      % Search direction: Generation II Feraligatr sprite beside ten character boxes, with O outside the limit.
      \lectureimage[2.22in]{slide8.jpg}

      \begin{keyidea}

        A storage decision can become a visible design constraint.

      \end{keyidea}

    \end{column}

  \end{columns}

  \vspace{0.05in}

  {\scriptsize Python's \texttt{str} type does not impose this ten-character limit; this is a general systems example.}

  \note{[Sources]\par
  Bulbapedia, ``Feraligatr (Pok\'emon),'' name origin and Generation II ten-character limit:\par
  \url{https://bulbapedia.bulbagarden.net/wiki/Feraligatr_(Pok\%C3\%A9mon)}.}

\end{frame}

\sectionframe{Part 2: Fundamental Python Types}

\begin{frame}{Integers Represent Whole Numbers Exactly}

  \begin{columns}[T,onlytextwidth]

    \begin{column}{0.48\textwidth}

      \begin{center}

        {\Large \texttt{12}}\\[0.12in]
        {\Large \texttt{-4}}\\[0.12in]
        {\Large \texttt{0}}\\[0.12in]
        {\Large \texttt{1\_000\_000}}

      \end{center}

    \end{column}

    \begin{column}{0.46\textwidth}

      \begin{itemize}

        \item Python names this type \texttt{int}.

        \item Integer literals contain no decimal point.

        \item Integers represent positive whole numbers, negative whole numbers, and zero.

        \item Integer arithmetic can produce another integer.

      \end{itemize}

    \end{column}

  \end{columns}

  \note{[Sources]\par
  Python documentation, Numeric Types:\par
  \url{https://docs.python.org/3/library/stdtypes.html\#numeric-types-int-float-complex}.}

\end{frame}

\begin{frame}{Python Integers Grow as More Digits Are Required}

  \begin{columns}[T,onlytextwidth]

    \begin{column}{0.52\textwidth}

      \begin{itemize}

        \item Many systems use fixed-width integers such as 8, 32, or 64 bits.

        \item Python integers use an arbitrary-precision representation.

        \item The representation grows to accommodate additional digits.

        \item The practical limit is available memory rather than a fixed 32-bit boundary.

      \end{itemize}

      \begin{exampleblockcp}

        CPython implements integers using an internal structure named \texttt{PyLongObject}.

      \end{exampleblockcp}

    \end{column}

    \begin{column}{0.42\textwidth}

      % Search direction: Python integer representation expanding as the number gains digits.
      \lectureimage[2.22in]{slide11.jpeg}

    \end{column}

  \end{columns}

  \note{[Sources]\par
  Python/C API documentation, Integer Objects:\par
  \url{https://docs.python.org/3/c-api/long.html}.}

\end{frame}

\begin{frame}{Floats Represent Numerical Values with Finite Precision}

  \begin{columns}[T,onlytextwidth]

    \begin{column}{0.48\textwidth}

      \begin{center}

        {\Large \texttt{3.5}}\\[0.13in]
        {\Large \texttt{-0.25}}\\[0.13in]
        {\Large \texttt{7.0}}

      \end{center}

      \begin{exampleblockcp}

        \texttt{7} and \texttt{7.0} describe the same numerical quantity but have different types.

      \end{exampleblockcp}

    \end{column}

    \begin{column}{0.46\textwidth}

      \begin{itemize}

        \item Python names this type \texttt{float}.

        \item CPython commonly uses 64-bit IEEE 754 binary floating point.

        \item A float provides roughly 15--17 significant decimal digits.

        \item Its numerical range is large but finite.

      \end{itemize}

    \end{column}

  \end{columns}

  \note{[Sources]\par
  Python documentation, Floating-Point Arithmetic:\par
  \url{https://docs.python.org/3/tutorial/floatingpoint.html}.}

\end{frame}

\begin{frame}{The Point ``Floats'' by Changing the Scale}

  \begin{columns}[T,onlytextwidth]

    \begin{column}{0.48\textwidth}

      \begin{itemize}

        \item The \textbf{sign} distinguishes positive and negative values.

        \item The \textbf{significand} contains the meaningful digits.

        \item The \textbf{exponent} determines the scale of the value.

        \item Changing the exponent effectively moves the binary point.

      \end{itemize}

      \begin{keyidea}

        Floating point supports both very small and very large magnitudes using a fixed amount of storage.

      \end{keyidea}

    \end{column}

    \begin{column}{0.46\textwidth}

      % Search direction: high-level floating-point diagram labeled sign, significand, and exponent.
      \lectureimage[2.25in]{slide13.png}

    \end{column}

  \end{columns}

\end{frame}

\begin{frame}{A Large Range Does Not Guarantee Exact Representation}

  \begin{columns}[T,onlytextwidth]

    \begin{column}{0.52\textwidth}

      \begin{itemize}

        \item A finite number of bits can represent only finitely many values.

        \item Many decimal fractions do not have an exact finite binary representation.

        \item Python therefore stores the nearest representable floating-point value.

      \end{itemize}

      \begin{exampleblockcp}

        \texttt{0.1 + 0.2} may display\\
        \texttt{0.30000000000000004}.

      \end{exampleblockcp}

    \end{column}

    \begin{column}{0.42\textwidth}

      % Search direction: decimal 0.1 approximated between nearby binary floating-point values.
      \lectureimage[2.18in]{slide14.jpg}

    \end{column}

  \end{columns}

  \note{[Sources]\par
  Python documentation, Representation Error:\par
  \url{https://docs.python.org/3/tutorial/floatingpoint.html\#representation-error}.}

\end{frame}

\begin{frame}{Strings Represent Unicode Text}

  \begin{columns}[T,onlytextwidth]

    \begin{column}{0.48\textwidth}

      \begin{center}

        {\Large \texttt{"hello"}}\\[0.12in]
        {\Large \texttt{"42"}}\\[0.12in]
        {\Large \texttt{"Cal Poly"}}\\[0.12in]
        {\Large \texttt{""}}

      \end{center}

    \end{column}

    \begin{column}{0.46\textwidth}

      \begin{itemize}

        \item Python names this type \texttt{str}.

        \item A string is an immutable sequence of Unicode characters.

        \item Its representation grows with the amount and kind of text.

        \item Quotes delimit a string literal; digits inside quotes remain text.

      \end{itemize}

    \end{column}

  \end{columns}

  \note{[Sources]\par
  Python documentation, Text Sequence Type --- \texttt{str}:\par
  \url{https://docs.python.org/3/library/stdtypes.html\#text-sequence-type-str}.}

\end{frame}

\begin{frame}{Boolean Values Represent Two Possibilities}

  \begin{columns}[T,onlytextwidth]

    \begin{column}{0.46\textwidth}

      \begin{center}

        {\Huge \texttt{True}}\\[0.28in]
        {\Huge \texttt{False}}

      \end{center}

    \end{column}

    \begin{column}{0.48\textwidth}

      \begin{itemize}

        \item Python names this type \texttt{bool}.

        \item The type contains exactly two literal values.

        \item Capitalization matters: \texttt{True} and \texttt{False}.

        \item Comparisons and Boolean operators will be introduced later.

      \end{itemize}

    \end{column}

  \end{columns}

  \note{[Sources]\par
  Python documentation, Boolean Type:\par
  \url{https://docs.python.org/3/library/stdtypes.html\#boolean-type-bool}.}

\end{frame}

\begin{frame}{\texttt{type()} Reports How Python Classifies a Value}

  \begin{center}

    \begin{tabular}{l l}
      \textbf{Expression} & \textbf{Result} \\
      \hline
      \texttt{type(7)} & \texttt{<class 'int'>} \\
      \texttt{type(7.0)} & \texttt{<class 'float'>} \\
      \texttt{type("7")} & \texttt{<class 'str'>} \\
      \texttt{type(True)} & \texttt{<class 'bool'>} \\
    \end{tabular}

  \end{center}

  \begin{keyidea}

    Python evaluates the argument first; \texttt{type()} then reports the resulting value's type.

  \end{keyidea}

\end{frame}

\sectionframe{Part 3: Operators Depend on Types}

\begin{frame}{Operand Types Select the Applicable Operator Rule}

  \begin{columns}[T,onlytextwidth]

    \begin{column}{0.49\textwidth}

      \begin{block}{Evaluation model}

        \centering
        operand values and types\\[0.06in]
        $\downarrow$\\[0.06in]
        operator compatibility rule\\[0.06in]
        $\downarrow$\\[0.06in]
        result value and type

      \end{block}

    \end{column}

    \begin{column}{0.45\textwidth}

      % Search direction: operator receiving typed operands and producing either a typed result or an error.
      \lectureimage[2.12in]{slide19.jpeg}

    \end{column}

  \end{columns}

  \begin{keyidea}

    An operator does not have one universal behavior independent of its operands.

  \end{keyidea}

\end{frame}

\begin{frame}{Numeric Types Support Arithmetic Operations}

  \begin{center}

    \begin{tabular}{l l l}
      \textbf{Expression} & \textbf{Value} & \textbf{Type} \\
      \hline
      \texttt{4 + 3} & \texttt{7} & \texttt{int} \\
      \texttt{4.0 + 3.0} & \texttt{7.0} & \texttt{float} \\
      \texttt{4 * 3} & \texttt{12} & \texttt{int} \\
      \texttt{4 * 3.0} & \texttt{12.0} & \texttt{float} \\
      \texttt{2 ** 3} & \texttt{8} & \texttt{int} \\
    \end{tabular}

  \end{center}

  \begin{exampleblockcp}

    The operator and operand types jointly determine the resulting value and type.

  \end{exampleblockcp}

\end{frame}

\begin{frame}{The Same \texttt{+} Symbol Can Mean Different Operations}

  \begin{center}

    \begin{tabular}{l l l}
      \textbf{Expression} & \textbf{Selected operation} & \textbf{Result} \\
      \hline
      \texttt{4 + 3} & integer addition & \texttt{7} \\
      \texttt{4.0 + 3.0} & floating-point addition & \texttt{7.0} \\
      \texttt{"4" + "3"} & string concatenation & \texttt{"43"} \\
    \end{tabular}

  \end{center}

  \vspace{0.10in}

\end{frame}

\begin{frame}{Strings Support Concatenation and Repetition}

  \begin{columns}[T,onlytextwidth]

    \begin{column}{0.48\textwidth}

      \begin{block}{Concatenation}

        \centering
        \texttt{"Cal" + " Poly"}\\[0.08in]
        $\rightarrow$\\[0.08in]
        \texttt{"Cal Poly"}

      \end{block}

    \end{column}

    \begin{column}{0.48\textwidth}

      \begin{block}{Repetition}

        \centering
        \texttt{"ha" * 3}\\[0.08in]
        $\rightarrow$\\[0.08in]
        \texttt{"hahaha"}

      \end{block}

    \end{column}

  \end{columns}

  \vspace{0.08in}

\end{frame}

\begin{frame}{Division Operators Can Produce Different Values and Types}

  \begin{center}

    \begin{tabular}{l l l}
      \textbf{Expression} & \textbf{Value} & \textbf{Type} \\
      \hline
      \texttt{4 / 2} & \texttt{2.0} & \texttt{float} \\
      \texttt{4 // 2} & \texttt{2} & \texttt{int} \\
      \texttt{5 / 2} & \texttt{2.5} & \texttt{float} \\
      \texttt{5 // 2} & \texttt{2} & \texttt{int} \\
      \texttt{5 // 2.0} & \texttt{2.0} & \texttt{float} \\
    \end{tabular}

  \end{center}

  \begin{keyidea}

    \texttt{/} always produces a \texttt{float}; floor division produces a \texttt{float} when a floating-point operand participates.

  \end{keyidea}

\end{frame}

\begin{frame}{Compatible Mixed Numeric Types Produce a Numeric Result}

  \begin{center}

    \begin{tabular}{l l l}
      \textbf{Expression} & \textbf{Value} & \textbf{Type} \\
      \hline
      \texttt{5 + 3.5} & \texttt{8.5} & \texttt{float} \\
      \texttt{4 * 2.0} & \texttt{8.0} & \texttt{float} \\
      \texttt{7 - 2.5} & \texttt{4.5} & \texttt{float} \\
      \texttt{7 // 2.0} & \texttt{3.0} & \texttt{float} \\
    \end{tabular}

  \end{center}

  \begin{exampleblockcp}

    When an arithmetic expression combines an \texttt{int} and a \texttt{float}, Python generally produces a \texttt{float}.

  \end{exampleblockcp}

\end{frame}

\begin{frame}{Unsupported Type Combinations Produce a \texttt{TypeError}}

  \begin{columns}[T,onlytextwidth]

    \begin{column}{0.48\textwidth}

      \begin{itemize}

        \item \texttt{4 + "3"}

        \item \texttt{"hello" - "h"}

        \item \texttt{"ha" * 2.5}

        \item \texttt{"4" / "2"}

      \end{itemize}

      \begin{cpdefinition}

        A \texttt{TypeError} indicates that an operation received a kind of value it does not support.

      \end{cpdefinition}

    \end{column}

    \begin{column}{0.46\textwidth}

    \end{column}

  \end{columns}

\end{frame}

\sectionframe{Part 4: Conversion, Prediction, and Explanation}

\begin{frame}{Explicit Conversion Requests a New Type}

  \begin{columns}[T,onlytextwidth]

    \begin{column}{0.48\textwidth}

      \begin{itemize}

        \item \texttt{int("12")} $\rightarrow$ \texttt{12}

        \item \texttt{float(8)} $\rightarrow$ \texttt{8.0}

        \item \texttt{str(3.5)} $\rightarrow$ \texttt{"3.5"}

      \end{itemize}

      \begin{keyidea}

        The conversion expression produces a new value of the requested type.

      \end{keyidea}

    \end{column}

    \begin{column}{0.46\textwidth}

      % Search direction: one value passing through int, float, and str conversion paths to new values.
      \lectureimage[2.18in]{slide27.jpg}

    \end{column}

  \end{columns}

\end{frame}

\begin{frame}{\texttt{int()} Produces an Integer Value}

  \begin{center}

    \begin{tabular}{l l l}
      \textbf{Expression} & \textbf{Value} & \textbf{Type} \\
      \hline
      \texttt{int("12")} & \texttt{12} & \texttt{int} \\
      \texttt{int(3.9)} & \texttt{3} & \texttt{int} \\
      \texttt{int(8.0)} & \texttt{8} & \texttt{int} \\
    \end{tabular}

  \end{center}

  \begin{exampleblockcp}

    Converting \texttt{3.9} does not round to the nearest integer; \texttt{int()} removes the fractional portion toward zero.

  \end{exampleblockcp}

\end{frame}

\begin{frame}{\texttt{float()} and \texttt{str()} Produce Different Representations}

  \begin{columns}[T,onlytextwidth]

    \begin{column}{0.47\textwidth}

      \begin{block}{Floating-point conversion}

        \texttt{float(8)} $\rightarrow$ \texttt{8.0}\\[0.09in]

        \texttt{float("3.5")} $\rightarrow$ \texttt{3.5}

      \end{block}

    \end{column}

    \begin{column}{0.47\textwidth}

      \begin{block}{String conversion}

        \texttt{str(3)} $\rightarrow$ \texttt{"3"}\\[0.09in]

        \texttt{str(3.5)} $\rightarrow$ \texttt{"3.5"}

      \end{block}

    \end{column}

  \end{columns}

  \begin{keyidea}

    Conversion changes how Python interprets the resulting value; it does not merely change how the value is displayed.

  \end{keyidea}

\end{frame}

\begin{frame}{Conversion Changes Which Operator Rule Applies}

  \begin{columns}[T,onlytextwidth]

    \begin{column}{0.47\textwidth}

      \begin{block}{Text remains text}

        \centering
        \texttt{"4" + "3"}\\[0.10in]
        $\rightarrow$\\[0.10in]
        \texttt{"43"}

      \end{block}

    \end{column}

    \begin{column}{0.47\textwidth}

      \begin{block}{Convert before addition}

        \centering
        \texttt{int("4") + int("3")}\\[0.10in]
        $\rightarrow$\\[0.10in]
        \texttt{7}

      \end{block}

    \end{column}

  \end{columns}

  \begin{exampleblockcp}

    The conversions change the operand types, so Python selects numerical addition instead of string concatenation.

  \end{exampleblockcp}

\end{frame}

\begin{frame}{Explicit and Implicit Numeric Conversion Are Different}

  \begin{center}

    \begin{tabular}{p{1.45in} p{2.50in} p{1.90in}}
      \textbf{Conversion} & \textbf{What happens} & \textbf{Example} \\
      \hline
      Explicit & The expression directly requests a target type. & \texttt{float(5)} \\
      Implicit numeric & Python applies a compatible mixed-numeric rule. & \texttt{5 + 3.5} \\
    \end{tabular}

  \end{center}

  \vspace{0.12in}

  \begin{columns}[T,onlytextwidth]

    \begin{column}{0.47\textwidth}

      \begin{exampleblockcp}

        \texttt{5 + 3.5} $\rightarrow$ \texttt{8.5}

      \end{exampleblockcp}

    \end{column}

    \begin{column}{0.47\textwidth}

      \begin{alertblock}{No guessing from text}

        \texttt{5 + "3.5"} $\rightarrow$ \texttt{TypeError}

      \end{alertblock}

    \end{column}

  \end{columns}

\end{frame}

\begin{frame}{Predict the Value and Type Before Running Python}

  \begin{columns}[T,onlytextwidth]

    \begin{column}{0.44\textwidth}

      \begin{enumerate}

        \item \texttt{7 + 2}

        \item \texttt{7 + 2.0}

        \item \texttt{7 / 2}

        \item \texttt{7 // 2}

        \item \texttt{"7" + "2"}

        \item \texttt{"ha" * 3}

        \item \texttt{int("7") + 2}

        \item \texttt{str(7) + "2"}

      \end{enumerate}

    \end{column}

    \begin{column}{0.50\textwidth}

      \begin{activity}

        For each expression, identify:

        \begin{enumerate}

          \item the operand types,

          \item the selected operation,

          \item the final value, and

          \item the final type.

        \end{enumerate}

      \end{activity}

    \end{column}

  \end{columns}

\end{frame}

\begin{frame}{Use Colab to Inspect Both the Value and Its Type}

  \begin{columns}[T,onlytextwidth]

    \begin{column}{0.47\textwidth}

      \begin{enumerate}

        \item Predict the value and type.

        \item Run the expression in one cell.

        \item Inspect the same expression using \texttt{type()}.

        \item Compare the observations with the prediction.

        \item Explain which type rule produced the result.

      \end{enumerate}

      \begin{exampleblockcp}

        \texttt{7 + 2.0}\\
        \texttt{type(7 + 2.0)}

      \end{exampleblockcp}

    \end{column}

    \begin{column}{0.47\textwidth}

    \end{column}

  \end{columns}

\end{frame}

\begin{frame}{Trace Types Through One Complete Expression}

  \begin{activity}

    Evaluate the expression without running Python first:

    \begin{center}

      {\LARGE \texttt{str(4 + 3.0) + " points"}}

    \end{center}

  \end{activity}

  \begin{enumerate}

    \item \texttt{4 + 3.0} produces \texttt{7.0}, a \texttt{float}.

    \item \texttt{str(7.0)} produces \texttt{"7.0"}, a \texttt{str}.

    \item String concatenation produces \texttt{"7.0 points"}.

  \end{enumerate}

  \begin{keyidea}

    To explain an expression, track both the value and the type produced by each subexpression.

  \end{keyidea}

\end{frame}

\end{document}
